% !TeX program = xelatex
\documentclass[12pt]{nextbook}

\UseTemplateSet{
  globalfonts = {libertinus,hindi},
  features = {hyperlinks}
}

\newcommand{\ReaderColumnSep}{4mm}
\newcommand{\ReaderVocabBaselineskip}{15pt}
\newcommand{\ReaderVocabSeparator}{}
\usepackage{paracol}
\usepackage{xcolor}
\usepackage{eso-pic}
\usepackage{needspace}

\vbadness=10000
\hbadness=10000
\hfuzz=3pt
\emergencystretch=1.5em

\makeatletter
\AtBeginDocument{%
  \@ifpackageloaded{hyperref}{%
    \pdfstringdefDisableCommands{%
      \def\ruby#1#2{#1}%
      \def\ReaderVocabText#1{#1}%
      \def\ReaderVocabTerm#1{#1}%
    }%
  }{}%
}
\makeatother

\providecommand{\ReaderColumnSep}{6mm}
\providecommand{\ReaderColumnRatio}{0.75}
\providecommand{\ReaderVocabBaselineskip}{20pt}
\providecommand{\ReaderVocabText}[1]{\textbf{#1}}
\providecommand{\ReaderVocabTerm}[1]{\textbf{#1}}
\providecommand{\ReaderVocabSeparator}{:}
\providecommand{\ReaderVocabItem}[2]{\noindent\ReaderVocabTerm{#1}\ReaderVocabSeparator\ #2\par}

\setlength{\columnsep}{\ReaderColumnSep}
\setlength{\columnseprule}{0pt}
\columnratio{\ReaderColumnRatio}
\swapcolumninevenpages

\newcounter{readersection}
\newwrite\vocabfile
\newif\iffirstreaderlabelwritten
\newlength{\readercolumnswidth}
\newlength{\readertitlegutter}
\setlength{\readercolumnswidth}{\dimexpr\textwidth-\columnsep\relax}
\setlength{\readertitlegutter}{.75\columnsep}

\makeatletter
\newcommand{\readerseparatorline}{%
  \if@mainmatter
    \AtTextLowerLeft{%
      \ifodd\value{page}%
        \put(\LenToUnit{\dimexpr .75\readercolumnswidth+.5\columnsep\relax},0){%
          \color{black!35}\rule{0.3pt}{\textheight}%
        }%
      \else
        \put(\LenToUnit{\dimexpr .25\readercolumnswidth+.5\columnsep\relax},0){%
          \color{black!35}\rule{0.3pt}{\textheight}%
        }%
      \fi
    }%
  \fi
}
\makeatother

\AddToShipoutPictureBG{\readerseparatorline}

\newcommand{\currentvocabfile}{\jobname-vocab-\thereadersection.voc}
\newcommand{\beginvocabwrite}{\immediate\openout\vocabfile=\currentvocabfile}
\newcommand{\finishvocabwrite}{\immediate\closeout\vocabfile}
\newcommand{\vocabitem}[2]{\ReaderVocabItem{#1}{#2}}
\newcommand{\printvocablist}{%
  \footnotesize
  \setlength{\parindent}{0pt}%
  \setlength{\parskip}{0pt}%
  \setlength{\baselineskip}{\ReaderVocabBaselineskip}%
  \raggedright
  \input{\currentvocabfile}%
}

\newcommand{\readersectiontitle}[1]{%
  \par
  \begingroup
  \ifodd\value{page}%
    \leftskip=\readertitlegutter%
    \rightskip=0.25\textwidth plus 1fil%
  \else
    \leftskip=\dimexpr .25\readercolumnswidth+.5\columnsep+\readertitlegutter\relax%
    \rightskip=0pt plus 1fil%
  \fi
  \section{#1}%
  \par
  \endgroup
}

\NewDocumentCommand{\voc}{omm}{%
  \ReaderVocabText{#2}%
  \IfNoValueTF{#1}{%
    \immediate\write\vocabfile{\string\vocabitem{#2}{#3}}%
  }{%
    \immediate\write\vocabfile{\string\vocabitem{#1}{#3}}%
  }%
}

\newenvironment{readersection}[1]{%
  \stepcounter{readersection}%
  \beginvocabwrite
  \Needspace{10\baselineskip}%
  \iffirstreaderlabelwritten\else
    \label{reader:firstpage}%
    \global\firstreaderlabelwrittentrue
  \fi
  \readersectiontitle{#1}%
  \begin{paracol}{2}
}{%
  \finishvocabwrite
  \switchcolumn
  \printvocablist
  \end{paracol}
}


\title{आरंभिक हिंदी पाठ्यपुस्तक}
\author{आकाश}
\date{\today}

\begin{document}
\frontmatter
\maketitle
\tableofcontents
\mainmatter
\chapter{प्रारंभिक पाठ}
\begin{readersection}{कमरे में पत्र}
\voc{यह}{this} \voc{एक}{a, one} \voc{कमरा}{room} है। कमरे \voc{में}{in} \voc{दो}{two} \voc[कुरसी]{कुरसियाँ}{chair} हैं। \voc{तीन}{three} \voc[पुस्तक]{पुस्तकें}{book}, \voc{क़लम}{pen}, \voc{अख़बार}{newspaper} \voc{और}{and} \voc{पत्र}{letter} \voc{मेज़}{table} \voc{पर}{on} हैं। \voc{यहाँ}{here} \voc{पानी}{water} \voc{नहीं}{not} है, \voc{लेकिन}{but} \voc{वहाँ}{there} \voc{घड़ा}{jar} है। घड़े में पानी है।

\voc{वह}{that} \voc{मकान}{house} है। मकान में एक \voc{लड़का}{boy}, एक \voc{लड़की}{girl}, एक \voc{आदमी}{man} और दो \voc[औरत]{औरतें}{woman} हैं। आदमी कमरे में है। लड़का कुरसी पर है। लड़की कमरे में है। औरतें मकान में हैं।

\voc{क्या}{what?} पत्र यहाँ है? नहीं, पत्र यहाँ नहीं है। पत्र मेज़ पर है। क़लम \voc{कहाँ}{where?} है? क़लम भी मेज़ पर है। पुस्तक कहाँ है? पुस्तक कुरसी पर नहीं है; पुस्तक मेज़ पर है।

मकान में एक \voc{चिड़िया}{bird} है। चिड़िया \voc[पिंजरा]{पिंजरे}{cage} में है। पिंजरा मेज़ पर नहीं है; पिंजरा कमरे में है। \voc[कुआँ]{कुएँ}{well} में पानी है, लेकिन घड़े में पानी नहीं है।

\voc{आगरा}{Agra} में एक लड़का है। \voc{कलकत्ता}{Calcutta} में दो लड़कियाँ हैं। लड़का वहाँ है, लड़कियाँ यहाँ नहीं हैं। कमरे में पुस्तकें हैं, मेज़ पर पत्र है, और पिंजरे में चिड़िया है।
\end{readersection}

\begin{readersection}{सीता का घर}
यह \voc{सीता}{Sita} \voc{का}{of} \voc{घर}{home} है। घर \voc{साफ़}{clean} है और घर की \voc{दीवार}{wall} \voc{सफ़ेद}{white} है। दीवार पर एक \voc{काला}{black} \voc{काग़ज़}{paper} है। काग़ज़ पर \voc{चार}{four} शब्द हैं। कमरे में मेज़ है, दो कुरसियाँ हैं, और मेज़ पर पुस्तक, \voc{किताब}{book}, क़लम और \voc{पेंसिल}{pencil} हैं।

घर का \voc{दरवाज़ा}{door} \voc{बड़ा}{big} है। दरवाज़े पर एक \voc{बूढ़ा}{old} आदमी \voc{खड़ा}{standing} है। कमरे में सीता का \voc{बेटा}{son}, उसकी \voc{बेटी}{daughter}, उसका \voc{भाई}{brother} और उसकी \voc{बहन}{sister} हैं। बेटा \voc{छोटा}{small} है, बेटी भी छोटी है। भाई बड़ा है, लेकिन बहन छोटी है। एक \voc{बच्चा}{child} कुरसी पर है और एक बच्चा मेज़ के पास है।

मेज़ के नीचे एक \voc{संदूक़}{box} है। संदूक़ में \voc{कपड़ा}{cloth} है। यह कपड़ा \voc{लाल}{red} है, वह कपड़ा सफ़ेद है। एक कपड़ा साफ़ है, लेकिन दूसरा कपड़ा \voc{मैला}{dirty} है। सीता कहती है, “यह कपड़ा \voc{बहुत}{very} मैला है, लेकिन वह लाल कपड़ा \voc{बढ़िया}{nice} है।”

घर के बाहर एक \voc{सुंदर}{beautiful} \voc{मंदिर}{temple} है। मंदिर बड़ा नहीं है; वह छोटा है। मंदिर के पास एक \voc{ऊँचा}{high} पेड़ है। पेड़ पर चिड़िया है। मंदिर के पीछे एक \voc{महल}{palace} है। महल बहुत बड़ा है और उसके सामने एक \voc{किला}{fort} है। किला पुराना है, लेकिन सुंदर है।

यह \voc{शहर}{city} \voc{दिल्ली}{Delhi} नहीं है। यह एक छोटा शहर है। शहर के पास एक \voc{गाँव}{village} है। गाँव में पानी है, कुआँ है, और दो छोटे मकान हैं। गाँव का कपड़ा \voc{सस्ता}{cheap} है। शहर का कपड़ा महँगा नहीं है, लेकिन गाँव का कपड़ा और सस्ता है।

\voc{मैं}{I} सीता की किताब देखता हूँ। किताब पर \voc{पाँच}{five} शब्द हैं। \voc[पाँचवाँ]{पाँचवाँ}{fifth} शब्द “\voc{देश}{country}” है। सीता कहती है, “मेरा देश बड़ा है। दिल्ली भारत की राजधानी है।” कमरे में सब लोग सुनते हैं। सीता का बेटा मेज़ पर पेंसिल रखता है, बेटी किताब खोलती है, और \voc{छोटा भाई}{younger brother} दरवाज़े पर खड़ा है।
\end{readersection}

\begin{readersection}{स्कूल में नमस्ते}
\voc{अभी}{now} मैं \voc{स्कूल}{school} में हूँ। स्कूल का कमरा साफ़ है। कमरे में मेज़, कुरसियाँ, पुस्तकें, किताबें, क़लम और पेंसिल हैं। मेज़ पर \voc[सब]{सब किताबें}{all} हैं। एक किताब \voc{भारत}{India} की है। भारत बड़ा देश है और दिल्ली भारत की \voc{राजधानी}{capital} है।

दरवाज़े पर एक आदमी है। वह स्कूल का \voc{मंत्री}{secretary} है। कमरे में सीता है, उसका \voc{पिता}{father} है, उसकी \voc[माता]{माँ}{mother} है, और उसके \voc{माता-पिता}{parents} दोनों यहाँ हैं। सीता की माँ कुरसी पर है। पिता मेज़ के पास हैं। मंत्री दरवाज़े पर \voc[बैठा]{बैठे}{sitting} नहीं हैं; वह खड़े हैं।

सीता का \voc{पति}{husband} भी स्कूल में है, और उसकी \voc{पत्नी}{wife} सीता है। पति कहता है, “\voc{नमस्ते}{greetings}, आप \voc{कैसे}{how} हैं?” सीता के पिता कहते हैं, “\voc{हाँ}{yes}, \voc{हम}{we} \voc{ठीक}{fine} हैं।” माँ कहती है, “मैं भी ठीक हूँ।”

कमरे में एक लड़का और एक लड़की हैं। लड़का कहता है, “मैं यहाँ हूँ।” लड़की कहती है, “तुम भी यहाँ हो।” मंत्री कहते हैं, “\voc{नमस्कार}{greetings}। आप सब स्कूल में हैं, और सब किताबें मेज़ पर हैं। आज \voc{काम}{work} है।” मेज़ पर काग़ज़ है। काग़ज़ पर छोटा काम है। यह काम कठिन नहीं है।

\voc{उस समय}{then} स्कूल में पानी नहीं है। घड़ा खाली है, लेकिन कुएँ में पानी है। पिता कहते हैं, “पानी वहाँ है।” लड़का दरवाज़े पर है, लड़की कमरे में है, और चिड़िया पिंजरे में है। सब लोग यहाँ हैं; काम भी यहाँ है।
\end{readersection}


\end{document}
